This chapter documents flight-controller bring-up on the SpeedyBee F405~V5 platform: Betaflight
telemetry experiments, sensor and motor commissioning on Betaflight, the later ArduPilot migration
for altitude-hold modes, and manual qualification with a conventional transmitter prior to wristband
demonstrations.

The aircraft carries a SpeedyBee F405~V5 autopilot, an OX32 55~A four-in-one ESC, four 1860~KV
motors, and a RadioMaster RX24T receiver fed by the ELRS link described in Section~\ref{subsec:elrs}.

%------------------------------------------------------------------------------
\subsection{Initial Betaflight Telemetry Scheme and Testing}
%------------------------------------------------------------------------------

Early work used Betaflight on the F405~V5 with the goal of closing an outer altitude loop on the
wristband using telemetry returned over ELRS\@.
Several CRSF telemetry frame types (GPS altitude, barometric altitude, attitude, battery) were
enabled and prioritized in Betaflight so altitude data could be forwarded to the ESP32 for display
and for hypothetical throttle corrections.
Firmware images were rebuilt after each scheduler change and flashed through the Betaflight
Configurator.

GPS altitude could be decoded but updated slowly and with coarse resolution.
Barometric altitude offered finer resolution in principle, yet under the tested ELRS schedules the
baro packets were not received consistently enough for closed-loop altitude control on the wearable.
Some frames appeared on the flight controller UART yet disappeared or were rate-limited across the
wireless hop.

These results showed that a wearable closed loop based on delayed or bursty telemetry was not
reliable for hover regulation, so the team abandoned outer-loop altitude trimming on the ESP32 in
favor of onboard ArduPilot altitude hold (see below) while still decoding barometric frames when available for
pilot awareness on the OLED.

\subsubsection{Betaflight Commissioning, Sensor Calibration, and Motor Verification}

Before first flight the following steps were performed over USB in Betaflight Configurator.

\textbf{Sensor calibration.}
Accelerometer calibration used the six-position wizard with the quad held still on each facet.
Gyro calibration used a motionless bench capture.

\textbf{Motor direction.}
Each motor was spun individually at minimum throttle on the Motors tab.
Two motors were reversed by swapping a pair of ESC phase leads until the prop-wash matched the
intended geometry.
Props were installed only after all four directions were verified.

\textbf{PID tuning policy.}
Factory Betaflight PID gains were retained; the team did not run a manual tuning loop because
the five-inch 1860~KV airframe responded adequately with defaults during hover shakedown.

\subsection{ArduPilot/ArduCopter Firmware Migration and AltHold Control}

To improve hover behavior, the flight controller was migrated from Betaflight to ArduPilot/
ArduCopter.
ArduPilot provides native \emph{AltHold} mode that closes altitude using the onboard barometer and
IMU, removing the need for low-rate telemetry on the wristband to close that loop.

The correct ArduCopter firmware target for the SpeedyBee F405~V5 was downloaded and flashed.
Mission Planner was then used to select the quad frame class, configure DShot (or the selected ESC
protocol), map the CRSF receiver to RC inputs, assign flight modes, and clear pre-arm checks.

With ArduPilot, the wristband continues to issue ordinary CRSF stick channels---it behaves like a
compact transmitter.
For example a ``hover'' intent can be realized by selecting AltHold via the mode channel while biasing
throttle near mid-stick; the inner-loop altitude regulator runs entirely on the autopilot.
Compared with the Betaflight telemetry experiment, this architecture trades wearable-side complexity
for deterministic stabilization and cleaner safety properties.

\subsection{Manual Flight Qualification}

A RadioMaster Pocket transmitter exercised the aircraft through motor checks, trim, and baseline
acro/angle flights before wristband-only demos.
The pilot (Chen Renang) adapted to handling and confirmed stable hover behavior, establishing a
reference for later comparisons against voice-command latency.
